\documentclass[english, parskip=full]{scrartcl}
\usepackage[utf8]{inputenc}  % Kodierung der Datei
\usepackage[english]{babel}  % Sprache auf Deutsch 
\usepackage{color}
\usepackage{graphicx, gensymb}
\usepackage{amsfonts, cancel, amssymb}
\usepackage{amsmath, amsthm, array, esint}
\usepackage{booktabs, multirow}  % schönere Tabellen
\usepackage{caption}  % Anpassung der Captions
\usepackage{scrlayer-scrpage}    % Manipulation des Seitenstils
\pagestyle{scrheadings}  % Stil für die Seite setzen
\automark{section}  % Abschnittsnamen als Seitenbeschriftung 
\setheadsepline{.5pt}    % Kopzeile durch Linie abtrennen
\usepackage[hidelinks]{hyperref}  % Links und weitere PDF-Features
\usepackage{typearea, tikz, pgffor, verbatim}
\usetikzlibrary{calc, quotes}
\usepackage[cache=false, outputdir=build]{minted}
\usepackage{placeins} % provides \FloatBarrier

\definecolor{bg}{rgb}{0.9,0.9,0.9}
\newcommand\numberthis{\addtocounter{equation}{1}\tag{\theequation}}
\newcommand{\bra}{}
\newcommand{\kom}{}
\newcommand{\brak}{}
\newcommand{\braket}{}
\newcommand{\intx}{}
\newcommand{\intr}{}
\newcommand{\kla}{}
\newcommand{\klam}{}
\newcommand{\klammer}{}
\newcommand{\oper}{}
\newcommand{\tecxt}{}
\newcommand{\Bra}{}
\newcommand{\Ket}{}
\renewcommand{\i}{\text{i}}
\newcommand{\e}{\text{e}}
\newcommand*{\Fourier}{\textsc{Fourier}\ }
\newcommand*{\F}{\mathcal{F}}
\newcommand*{\sinc}{\text{sinc\,}}
\newcommand{\conv}{\mathop{\scalebox{3.5}{\raisebox{-0.38ex}{\makebox[0.5\width][c]{$\ast$}}}}\limits}

\newcommand*{\titel}{02 Berry phase and Berry connection}
\newcommand*{\autor}{Joachim Schwardt}

\hypersetup{pdfauthor={\autor}, pdftitle={\titel}}

%\titlehead{titlehead}
\subject{Topological Insulators and Superconductors}
\title{\titel}
\author{\autor}
\date{\today}

\begin{document}

%\begin{titlepage}
\maketitle  % Titel setzen
%\tableofcontents  % Inhaltsverzeichnis setzen
%\end{titlepage}

\section*{Problem 2.1: Berry phase for a four-spin loop}
Consider a path through the four spinor states
\begin{align}
\Ket | u_0 \rangle &= \begin{pmatrix}
1 \\ 0
\end{pmatrix} &&& \Ket | u_1 \rangle &= \frac{1}{\sqrt{2}}\begin{pmatrix}
1 \\ 1
\end{pmatrix} &&& \Ket | u_2 \rangle &= \begin{pmatrix}
0 \\ 1
\end{pmatrix} &&& \Ket | u_3 \rangle &= \frac{1}{\sqrt{2}}\begin{pmatrix}
1 \\ \i
\end{pmatrix} .
\end{align}
We know that, in the discrete setting, the Berry phase is simply related to the quantity
\begin{align}
S &:= \prod_{j\in\tecxt\text{loop}} \brak\langle u_j | u_{j+1} \rangle = \brak\langle u_0 | u_1 \rangle\brak\langle u_1 | u_2 \rangle \brak\langle u_2 | u_3 \rangle \brak\langle u_3 | u_0 \rangle = \label{eqn:S discrete Berry phase}\\
&= \frac{1}{\sqrt{2}}\cdot \frac{1}{\sqrt{2}} \cdot \frac{\i}{\sqrt{2}} \cdot \frac{1}{\sqrt{2}} = \frac{\i}{4}.
\end{align}
The Berry phase is then given by $\phi := -\tecxt\text{Im\,}\log S$.
Note that in general
\begin{align}
\log z &= \log r\e^{\i\chi} = \log r + \i\chi + 2\pi\i n,
\end{align}
where $n\in\mathbb{Z}$ may be any integer.
However, since the Berry phase is only defined modulo $2\pi$, we may write
\begin{align}
\phi &= -\tecxt\text{Im\,}\log S = -\tecxt\text{Im\,} |S| -\tecxt\text{Im\,}\i\arg S = -\arg S.
\end{align}
Since in our case $\arg S = \frac{\pi}{2}$, we find $\phi = -\frac{\pi}{2}$.
%\begin{align}
%\phi &= -\tecxt\text{Im\,}\log \frac{\i}{4} = -\tecxt\text{Im\,}\log \i = -\tecxt\text{Im\,}\i\frac{\pi}{2} = -\frac{\pi}{2}.
%\end{align}

\newpage
\section*{Problem 2.2: Two-dimensional massive Dirac Hamiltonian}
Consider the two-dimensional massive Dirac Hamiltonian
\begin{align}
H(k_x,k_y) &= m\tau_z + k_x\tau_x + k_y\tau_y.
\end{align}
Since $(\vec{A}\cdot\vec{\tau})^2 = \vec{A}\,^2$, we find the eigenvalues 
\begin{align}
\epsilon_\pm &= \pm\sqrt{H^2} = \pm\sqrt{m^2 + \vec{k}^2},
\end{align}
where $\vec{k} := (k_x, k_y)^\top$.
The eigenstates $\Ket | \pm \rangle := (a, b)^\top$ are found by solving the system of equations
\begin{align}
H\Ket | \pm \rangle &= \begin{pmatrix}
ma + (k_x-\i k_y)b \\ (k_x + \i k_y)a - mb
\end{pmatrix} \overset{!}{=} \epsilon_\pm\Ket | \pm \rangle = \epsilon_\pm \begin{pmatrix}
a \\ b
\end{pmatrix}.
\end{align}
The second row yields
\begin{align}
b &= \frac{k_x + \i k_y}{m + \epsilon_\pm} a,
\end{align}
which leads to the eigenstates 
\begin{align}
\Ket | \pm \rangle &= N_\pm \begin{pmatrix}
m + \epsilon_\pm \\ k_x + \i k_y
\end{pmatrix}.
\end{align}
The normalization constant is
\begin{align}
N_\pm^2 &= (m+\epsilon_\pm)^2 + (k_x^2 + k_y^2) = 2\epsilon_\pm\kla\left( m + \epsilon_\pm \right).
\end{align}
Thus the eigenstates can be expressed as
\begin{align}
\Ket | \pm \rangle &= \frac{1}{\sqrt{2\epsilon (m \pm \epsilon)}} \begin{pmatrix}
m \pm \epsilon \\ k_x + \i k_y
\end{pmatrix},
\end{align}
where $\epsilon := |\epsilon_\pm| = \sqrt{m^2 + \vec{k}^2}$.
The Berry connection is given by
\begin{align}
\vec{A} &:= \braket\langle \pm | \i\partial_{\vec{k}} | \pm \rangle,
\end{align}
where we may interpret our two-dimensional setting as one of complex variables.
Then we can express the Berry connection as
% the components are
%\begin{align}
%A_1 &= \frac{1}{\epsilon (m\pm \epsilon)} \begin{pmatrix}
%m\pm\epsilon & k_x -\i k_y
%\end{pmatrix} \i\partial_{k_x} \begin{pmatrix}
%m\pm\epsilon \\ k_x +\i k_y
%\end{pmatrix} = \frac{k_y + \i k_x}{\epsilon (m\pm \epsilon)}, \\
%A_2 &= \frac{1}{\epsilon (m\pm \epsilon)} \begin{pmatrix}
%m\pm\epsilon & k_x -\i k_y
%\end{pmatrix} \i\partial_{k_y} \begin{pmatrix}
%m\pm\epsilon \\ k_x +\i k_y
%\end{pmatrix} = \frac{\i k_y -  k_x}{\epsilon (m\pm \epsilon)}.
%\end{align}
\begin{align}
A(k) &= \braket\langle \pm | \i\partial_k | \pm \rangle = \frac{1}{2\epsilon (m\pm \epsilon)} \begin{pmatrix}
m\pm\epsilon & k^\ast
\end{pmatrix} \i\partial_{k} \begin{pmatrix}
m\pm\epsilon \\ k
\end{pmatrix} = \\
&= \frac{\i }{2\epsilon (m\pm \epsilon)}\kla\left( (m\pm \epsilon)\frac{\pm k}{\epsilon} +  k^\ast \right) = \frac{\pm\i k}{2\epsilon^2} + \frac{\i k^\ast}{2\epsilon (m\pm\epsilon)}.
\end{align}
Let $k(t)=k\e^{\i t}$ for $t\in[0,1)$ be a circular loop with radius $k$ in the parameter space. 
Then we may compute the Berry phase via
\begin{align}
\phi &= \oint A(k)\,\tecxt\text{d}k = \intx\int_{0}^{2\pi} A(t) \,\i k(t) \mathrm{d}t = \\
&= \intx\int_{0}^{2\pi}\frac{\mp k^2}{2\epsilon^2} \underbrace{\e^{2\i t}}_{\rightarrow\, 0}\, \mathrm{d}t +  \frac{-1}{2\epsilon (m\pm \epsilon)}\intx\int_{0}^{2\pi} |k|^2\, \mathrm{d}t = \\
&= -\frac{\pi |k|^2 (m\mp \epsilon)}{\epsilon (m^2 - \epsilon^2)} = \frac{\pi (m\mp \epsilon)}{\epsilon} = \pi \kla\left( \frac{m}{\sqrt{m^2 + |k|^2}} \mp 1 \right).
\end{align}
Now we will use the second approach to compute the same result.
Note that the Berry curvature is here given by a pseudo-cross product, i.e. only the $k_z$-component 
\begin{align}
\Omega(\vec{k}) &= |\partial_{\vec{k}}\times \vec{A}(\vec{k})| = |\partial_{k_x} A_y - \partial_{k_y} A_x |.
\end{align}
We will first compute the partial derivatives of our eigenstates
\begin{align}
\partial_{k_x}\Ket | \pm \rangle &= \partial_{k_x}\frac{1}{\sqrt{2\epsilon(m\pm\epsilon)}} \begin{pmatrix}
m\pm \epsilon \\ k_x + \i k_y 
\end{pmatrix} = \\
&= -\frac{\partial_{k_x}\epsilon}{2\epsilon}\Ket | \pm \rangle - \frac{\partial_{k_x}(m\pm \epsilon)}{2(m\pm \epsilon)}\Ket | \pm \rangle + \frac{1}{\sqrt{2\epsilon(m\pm\epsilon)}} \begin{pmatrix}
0 \\ 1
\end{pmatrix} = \\
&= -\frac{1}{2}\kla\left( \frac{k_x}{\epsilon^2} \pm \frac{k_x}{\epsilon(m\pm \epsilon)} \right)\Ket | \pm \rangle + \frac{1}{\sqrt{2\epsilon(m\pm\epsilon)}} \begin{pmatrix}
0 \\ 1
\end{pmatrix},\\
\partial_{k_y}\Ket | \pm \rangle &= -\frac{1}{2}\kla\left( \frac{k_y}{\epsilon^2} \pm \frac{k_y}{\epsilon(m\pm \epsilon)} \right)\Ket | \pm \rangle + \frac{1}{\sqrt{2\epsilon(m\pm\epsilon)}} \begin{pmatrix}
0 \\ \i
\end{pmatrix}.
\end{align}
Then the Berry curvature can be obtained from
\begin{align}
\Omega &= 2\tecxt\text{Im\,}\brak\langle \partial_{k_x}\pm | \partial_{k_y}\pm \rangle = \\
&= \frac{1}{\epsilon (m\pm\epsilon)} - \frac{k_x^2+k_y^2}{\epsilon^2}\kla\left( 1\pm \frac{\epsilon}{m\pm\epsilon} \right)\frac{1}{\sqrt{2\epsilon(m\pm\epsilon)}^2} = \\
&= -\frac{k^2}{2\epsilon^3(m\pm\epsilon)}\kla\left( 1 \pm \frac{\epsilon}{m\pm\epsilon} - \frac{2\epsilon^2}{k^2} \right) = \\
&= -\frac{1}{2\epsilon^3(m\pm \epsilon)}\kla\left( k^2 - 2\epsilon^2 \pm \frac{\epsilon (\epsilon^2 - m^2)}{m\pm\epsilon} \right) = \\
&= \frac{1}{2\epsilon^3(m\pm \epsilon)^2} \kla\left( (\epsilon^2 + m^2)(m\pm\epsilon) \mp\epsilon(\epsilon^2-m^2) \right) = \\
&= \frac{1}{2\epsilon^3(m\pm \epsilon)^2} \kla\left( m^3\pm 2m^2\epsilon + m\epsilon^2 \right) = \frac{m}{2\epsilon^3} .
\end{align}
This is just the $z$-component of the general formula $\vec{\Omega} = \frac{\vec{d}}{2d^3}$, which would have been a much faster approach.
%It is more convenient to express $A(k)$ as a vector here, i.e. 
%\begin{align}
%\vec{A} &= \frac{\i}{2\epsilon^2(m\pm \epsilon)}\begin{pmatrix}
%\pm k_x(m\pm\epsilon) + k_x\epsilon \\ 
%\pm k_y(m\pm\epsilon) - k_y\epsilon
%\end{pmatrix} = \frac{\i}{2\epsilon^2(m\pm \epsilon)}\begin{pmatrix}
%\pm k_x m + 2k_x\epsilon \\ 
%\pm k_ym
%\end{pmatrix} = \\
%&= \frac{\pm\i}{2\epsilon^2}\begin{pmatrix}
% k_x \\ 0
%\end{pmatrix} + \frac{\i}{2\epsilon^2(\epsilon\pm m)} \begin{pmatrix}
%\pm k_x\epsilon \\ 
%k_ym
%\end{pmatrix}.
%\end{align}
%Then using $\epsilon^2 = m^2+k_x^2+k_y^2$ the Berry curvature becomes
%\begin{align}
%\Omega &= \left\vert \partial_{k_x} \frac{ k_ym}{2\epsilon^2( \epsilon \pm m)} - \partial_{k_y} \frac{\pm k_x}{2\epsilon^2} - \partial_{k_y} \frac{\pm k_x\epsilon}{2\epsilon^2( \epsilon \pm m)} \right\vert = \\
%&= \left\vert \frac{\pm k_xk_y}{\epsilon^4} - \frac{k_xk_y (m\mp\epsilon)}{\epsilon^4(\epsilon\pm m)} + \frac{k_xk_y (-m\mp \epsilon)}{2\epsilon^3(\epsilon\pm m)^2} \mp \frac{k_xk_y}{2\epsilon^3(\epsilon\pm m)} \right\vert = \\
%&= \frac{k_xk_y}{2\epsilon^4}\left\vert 2 + \frac{2(\epsilon^2 - m^2)}{(\epsilon\pm m)^2} - \frac{\epsilon}{\epsilon\pm m} - \frac{\epsilon}{\epsilon\pm m} \right\vert = \\
%&= \frac{k_xk_y}{\epsilon^4(\epsilon\pm m)^2}\left\vert (\epsilon\pm m)^2 + \epsilon^2 - m^2 - \epsilon (\epsilon\pm m) \right\vert = \\
%&= \frac{k_xk_y}{2\epsilon^4(\epsilon\pm m)^2}\left\vert \epsilon^2 \pm \epsilon m \right\vert = \\
%&= \frac{m}{2\epsilon^3}  \left\vert \frac{k_xk_y}{m(\epsilon\pm m)} \right\vert.
%\end{align}
%The literature gives this result without the complicated absolute value factor.
Then the Berry phase is given by an integral over the surface, corresponding to the boundary $k(t)$ used above.
The integral can most easily be evaluated in polar $k$-coordinates, since there is no angular dependence.
This yields
\begin{align}
\phi &= \intx\int_{ }^{ }\Omega(k)\, 2\pi k\mathrm{d}k = 2\pi\intx\int_{0}^{k} \frac{m}{2\epsilon(k')^3}\, k'\mathrm{d}k' = \pi m\intx\int_{0}^{k}\frac{k'}{\kla\left( m^2 + k'^2 \right)^{3/2}}\, \mathrm{d}k' = \\
&= -\pi m\kla\left( m^2 + k'^2 \right)^{-1/2}\Big\vert_0^k  = \pi \kla\left( 1 - \frac{m}{\epsilon} \right).
\end{align}
Note that this is not quite equivalent to the previous expression, so there must be some mistake along the way, but there is most likely just some error/confusion with the $\pm$ signs.

%We may again reinterpret this as a complex derivative, since $\partial_z f(z) = \partial_x\tecxt\text{Re\,}f + \i\partial_y\tecxt\text{Im\,}f$.
%Note that $A_x\equiv \tecxt\text{Re\,}A$ and $A_y \equiv \tecxt\text{Im\,}A$.
%Using the Cauchy-Riemann equations $\partial_x\tecxt\text{Re\,}f = \partial_y\tecxt\text{Im\,}f$ and $\partial_y\tecxt\text{Re\,}f = -\partial_x\tecxt\text{Im\,}f$ we find
%\begin{align}
%\Omega (k) &\equiv \partial_{k_x}\tecxt\text{Im\,}A(k) - \partial_{k_y}\tecxt\text{Re\,}A(k) = \partial_{k_x}\tecxt\text{Im\,}A(k) + \partial_{k_y}\tecxt\text{Im\,}A(k) \equiv \left\vert\partial_k A(k)\right\vert.
%\end{align}



\newpage
\section*{Problem 2.3: Berry phase in a magnetic field}
Consider a spin-1/2 particle in a uniform magnetic field $\vec{B} = B\vec{n}$. 
Its Hamiltonian is given by
\begin{align}
H &= -\gamma\vec{B}\cdot \vec{S} = -\frac{\gamma\hbar B}{2}\vec{n}\cdot\vec{\tau}.
\end{align}
Note that $\vec{n}  =(\cos\phi\sin\theta, \sin\phi\sin\theta, \cos\theta)^\top$, which lets us express the Hamiltonian in spherical coordinates
\begin{align}
H &= -\frac{\gamma\hbar B}{2} \kla\left( \cos\phi\sin\theta\tau_x + \sin\phi\sin\theta\tau_y + \cos\theta\tau_z \right).
\end{align}
Let $\Ket | \uparrow_{\vec{n}} \rangle = (a, b)^\top$ be the eigenstates of $H$ with eigenvalue $\lambda_{\vec{n}}$.
Note that the eigenvalues are simply found by
\begin{align}
H^2 &= \kla\left( \frac{\gamma\hbar B}{2} \right)^2 \underbrace{(\vec{n}\cdot\vec{\tau})^2}_{=\,\vec{n}\,^2\,=\, 1} = \lambda_{\vec{n}}^2 &&\Rightarrow & \lambda_{\vec{n}} &= \pm \frac{\gamma\hbar B}{2}.
\end{align}
The spin-up state can thus be found by solving the system of equation given by 
\begin{align}
H\Ket | \uparrow_{\vec{n}} \rangle &=  -\frac{\gamma\hbar B}{2} \begin{pmatrix}
\cos\theta a + \sin\theta \e^{-\i\phi}b \\
\sin\theta\e^{\i\phi} a - \cos\theta b
\end{pmatrix} \overset{!}{=} \lambda_{\vec{n}}\Ket | \uparrow_{\vec{n}} \rangle = \frac{\gamma\hbar B}{2}\begin{pmatrix}
a \\ b
\end{pmatrix}.
\end{align}
The second row gives us the relation 
\begin{align}
b =  a \e^{\i\phi} \frac{\sin\theta}{1 + \cos\theta}  =  a \e^{\i\phi}  \frac{2\sin\frac{\theta}{2}\cos\frac{\theta}{2}}{ 2\cos^2\frac{\theta}{2}} =  a \e^{\i\phi} \tan\frac{\theta}{2}.
\end{align}
%for $\lambda_{\vec{n}} = -\frac{\gamma\hbar B}{2}$ the second row gives us 
%\begin{align}
%b =  a \e^{\i\phi} \frac{\sin\theta}{-1 + \cos\theta}  =  a \e^{\i\phi}  \frac{2\sin\frac{\theta}{2}\cos\frac{\theta}{2}}{ -2\sin^2\frac{\theta}{2}} =  -a \e^{\i\phi} \cot\frac{\theta}{2}.
%\end{align}
%Inserting into the first row yields
%\begin{align}
%\cos\theta + \sin\theta\tan\theta &= \lambda_{\vec{n}} &&\Rightarrow & 1 &= \lambda_{\vec{n}}\cos\theta.
%\end{align}
%Thus we obtain the eigenvalues as $\lambda_{\vec{n}} = \frac{1}{\cos\theta}$.
Thus we can write our eigenstates as
\begin{align}
\Ket | \uparrow_{\vec{n}} \rangle &\propto \begin{pmatrix}
1 \\ \e^{\i\phi}\tan\theta/2
\end{pmatrix},
\end{align}
where the expression currently has a squared magnitude of
\begin{align}
N_+^2 &= 1 + \tan^2\frac{\theta}{2} = \frac{1}{\cos^2\frac{\theta}{2}} .
\end{align}
%respectively. && \tecxt\text{and} & N_-^2 &= 1 + \cot^2\frac{\theta}{2} = \frac{1}{\sin^2\frac{\theta}{2}}
Normalizing then yields
\begin{align}
\Ket | \uparrow_{\vec{n}} \rangle &= \begin{pmatrix}
\cos\frac{\theta}{2} \\ \e^{\i\phi}\sin\frac{\theta}{2}
\end{pmatrix}.
\end{align}
Next, consider $N$ states $\Ket | \uparrow_j \rangle := \Ket | \uparrow_{\vec{n}(\phi_j)} \rangle$ with equidistant values of $\phi_j := \frac{2\pi j}{N}$ for $j = 0,\dots,N-1$ for a fixed $\theta$.
We want to compute the Berry phase, and for that matter once again start with the product formula \eqref{eqn:S discrete Berry phase}
\begin{align}
S &= \prod_{j=0}^{N-1} \brak\langle \uparrow_j | \uparrow_{j+1} \rangle =  \prod_{j=0}^{N-1} \cos^2\frac{\theta}{2} + \kla\left( \e^{-\i\phi_j}\e^{\i\phi_{j+1}} \right) \sin^2\frac{\theta}{2} = \\
&= \prod_{j=0}^{N-1} \cos^2\frac{\theta}{2} + \e^{\i \frac{2\pi}{N}} \sin^2\frac{\theta}{2} = \kla\left( \cos^2\frac{\theta}{2} + \e^{\i \frac{2\pi}{N}} \sin^2\frac{\theta}{2} \right)^N.
\end{align}
Now the berry phase is given by
\begin{align}
\phi &= -\tecxt\text{Im\,}\log S = -N\tecxt\text{Im\,} \log\kla\left( \cos^2\frac{\theta}{2} + \e^{\i \frac{2\pi }{N}} \sin^2\frac{\theta}{2} \right) = \\
&= -N \arg\kla\left( \cos^2\frac{\theta}{2} + \cos\frac{2\pi }{N} \sin^2\frac{\theta}{2}  + \i\sin\frac{2\pi}{N}\sin^2\frac{\theta}{2} \right) = \\
&= -N \arctan\kla\left( \frac{\sin\frac{2\pi }{N}\sin^2\frac{\theta}{2}}{\cos^2\frac{\theta}{2} + \cos\frac{2\pi }{N}\sin^2\frac{\theta}{2}} \right).
\end{align}
In the limit $N\rightarrow\infty$ we find $\cos\frac{2\pi}{N}\rightarrow 1 + \mathcal{O}\kla\left( \frac{1}{N^2} \right)$ and $\sin\frac{2\pi}{N}\rightarrow \frac{2\pi}{N} + \mathcal{O}\kla\left( \frac{1}{N^3} \right)$.
Thus, using $\arctan(x) \rightarrow x + \mathcal{O}(x^3)$, we may write (not completely rigorously)
\begin{align}
\lim_{N\rightarrow\infty} \phi &= -\lim_{N\rightarrow\infty}N\arctan\kla\left( \frac{\frac{2\pi }{N}\sin^2\frac{\theta}{2} + \mathcal{O}\kla\left( \frac{1}{N^3} \right)}{\cos^2\frac{\theta}{2} + \sin^2\frac{\theta}{2} + \mathcal{O}\kla\left( \frac{1}{N^2} \right)} \right) = \\
&= -\lim_{N\rightarrow\infty} N\arctan\kla\left( \frac{2\pi}{N}\sin^2\frac{\theta}{2} \right) = \\
&= -\lim_{N\rightarrow\infty} N \frac{2\pi}{N}\sin^2\frac{\theta}{2} + \mathcal{O}\kla\left( \frac{1}{N^2} \right) = -2\pi\sin^2\frac{\theta}{2}.
\end{align}

Now consider the continuous case.
We have two parameters, denoted as $R_1 := \theta$ and $R_2 := \phi$.
Also let $\vec{R} := (R_1,R_2)^\top$.
Then the Berry connection/potential $\vec{A} = (A_1, A_2)^\top$ is defined by
\begin{align}
A_1 &:= \braket\langle \uparrow_{\vec{n}} | \i\partial_{R_1} | \uparrow_{\vec{n}} \rangle = \begin{pmatrix}
\cos\frac{\theta}{2} & \e^{-\i\phi}\sin\frac{\theta}{2}
\end{pmatrix} \i\partial_\theta\begin{pmatrix}
\cos\frac{\theta}{2} \\ \e^{\i\phi}\sin\frac{\theta}{2}
\end{pmatrix} = \\
&= \i\kla\left( -\cos\frac{\theta}{2}\sin\frac{\theta}{2} + \sin\frac{\theta}{2}\cos\frac{\theta}{2} \right) = 0, \\
A_2 &:= \braket\langle \uparrow_{\vec{n}} | \i\partial_{R_2} | \uparrow_{\vec{n}} \rangle = \begin{pmatrix}
\cos\frac{\theta}{2} & \e^{-\i\phi}\sin\frac{\theta}{2}
\end{pmatrix} \i\partial_\phi\begin{pmatrix}
\cos\frac{\theta}{2} \\ \e^{\i\phi}\sin\frac{\theta}{2}
\end{pmatrix} = \\
&= \i\kla\left( 0 + \i\sin^2\frac{\theta}{2}  \right) = -\sin^2\frac{\theta}{2}.
\end{align}
Now let $t\in [0,2\pi)$, i.e. $\tecxt\text{d}\vec{R} = (0, \tecxt\text{d}t)^\top$ (since $\theta$ is fixed, we have $\tecxt\text{d}\theta = 0$).
Then the Berry phase for this particular loop is given by
\begin{align}
\phi &= \int_{0}^{2\pi} \vec{A} \cdot \mathrm{d}\vec{R} = \int_{0}^{2\pi} -\sin^2\frac{\theta}{2}\, \mathrm{d}t = -2\pi\sin^2\frac{\theta}{2},
\end{align}
which is exactly what we found in the limit of our discrete solution.
The berry curvature is given by
\begin{align}
\Omega &= \partial_{R_1}A_2 - \partial_{R_2}A_1 = \partial_\theta \kla\left( -\sin^2\frac{\theta}{2} \right) = -2\sin\frac{\theta}{2}\cos\frac{\theta}{2} \cdot \frac{1}{2} = -\frac{1}{2}\sin\theta.
\end{align}

\newpage
\section*{Problem 2.4: The anisotropic QWZ model}
The lattice Hamiltonian of the anisotropic Qi-Wu-Zhang (QWZ) model is 
%\begin{align*}
%H &= -\eta\sum_{j=1}^{N-1}\sum_{k=1}^N\frac{\tau_{\alpha\beta}^z + \i\tau_{\alpha\beta}^x}{2}c_{j+1,k,\alpha}^\dagger c_{j,k,\beta} + \tecxt\text{h.c.} - \\
%&\hspace{0.5cm} - \sum_{j=1}^{N-1}\sum_{k=1}^N\frac{\tau_{\alpha\beta}^z + \i\tau_{\alpha\beta}^x}{2}c_{j,k+1,\alpha}^\dagger c_{j,k,\beta} + \tecxt\text{h.c.} + \\
%&\hspace{0.5cm} + u\sum_{j=1}^{N}\sum_{k=1}^N\tau_{\alpha\beta}^zc_{j,k,\alpha}^\dagger c_{j,k,\beta} \numberthis.
%\end{align*}
\begin{align*}
H &= -\eta\sum_{m_x=1}^{N-1}\sum_{m_y=1}^N\oper | m_x+1,m_y \rangle\langle m_x,m_y |\otimes\frac{\tau_z + \i\tau_x}{2} + \tecxt\text{h.c.} - \\
&\hspace{0.5cm} - \sum_{m_x=1}^{N-1}\sum_{m_y=1}^N\oper | m_x,m_y+1 \rangle\langle m_x,m_y |\otimes\frac{\tau_z + \i\tau_y}{2} + \tecxt\text{h.c.} + \\
&\hspace{0.5cm} + u\sum_{j=1}^{N}\sum_{k=1}^N\oper | m_x,m_y \rangle\langle m_x,m_y |\otimes\tau_z \numberthis.
\end{align*}
We can introduce momentum operators to express the annihilation operators
\begin{align}
\Ket | m_x,m_y \rangle &= \frac{1}{\sqrt{N}}\sum_{k_x,k_y} \e^{\i k_x a m_x + \i k_y am_y}\Ket | k_x,k_y \rangle.
\end{align}
Inserting this into the Hamiltonian yields
\begin{align*}
H &= -\frac{1}{N}\sum_{k_x,k_x',k_y,k_y',m_x,m_y}\e^{\i (k_x-k_x') a m_x + \i (k_y-k_y') am_y}\oper | k_x,k_y \rangle\langle k_x',k_y' |\otimes\\
&\hspace{0.7cm} \otimes \kla\left( \eta\e^{\i k_xa}\frac{\tau_z + \i\tau_x}{2} + \e^{\i k_ya}\frac{\tau_z + \i\tau_y}{2} + \tecxt\text{h.c.} + u\tau_z\right) =  \numberthis\\
&= -\sum_{k_x,k_y} \oper | k_x,k_y \rangle\langle k_x,k_y | \otimes\\
&\hspace{0.7cm}  \otimes \kla\Big( \tau_z\klam\left[ u + \eta\cos(k_xa) + \cos(k_ya) \right] + \eta\tau_x\sin(k_xa) + \tau_y\sin(k_ya) \Big). \numberthis
\end{align*}
We find the Hamiltonians 
\begin{align}
H(k_x,k_y) &= \braket\langle k_x,k_y | H | k_x,k_y \rangle = \\
&= \tau_z\klam\left[ u + \eta\cos(k_xa) + \cos(k_ya) \right] + \eta\tau_x\sin(k_xa) + \tau_y\sin(k_ya) =: \vec{d}\cdot\vec{\tau},
\end{align}
which lead to the energy bands
\begin{align}
E(k_x,k_y) &= \pm |\vec{d}| = \pm \sqrt{\eta^2\sin^2(k_xa) + \sin^2(k_ya) + \klam\left[ u + \eta\cos(k_xa) + \cos(k_ya) \right]^2}.
\end{align}
We know that the Chern number can only change, when $\vec{d}=0$.
This leads to the equations
\begin{align}
0 &= \begin{pmatrix}
\eta\sin(k_xa) \\ \sin(k_ya) \\ u + \eta\cos(k_xa) + \cos(k_ya)
\end{pmatrix},
\end{align}
which has several solutions listed in table \ref{table:chern d vec zero crossings}.
\begin{table}[!ht]
\centering
\begin{tabular}{c|c||c}
$k_xa$ & $k_ya$ & $u(\eta)$ \\ \hline
0 & 0 & $-1-\eta$ \\ \hline
$\pi$ & 0 & $-1+\eta$ \\ \hline
0 & $\pi$ & $1-\eta$ \\ \hline
$\pi$ & $\pi$ & $1+\eta$ \\
\end{tabular}
\caption{Momenta corresponding to zero-crossings of $\vec{d}(k_x,k_y)$ lead to different values of the parameter $u$ as a function of $\eta$.}
\label{table:chern d vec zero crossings}
\end{table}
\\
In the resulting $u,\eta$ parameter space, we find that there may in principle be
\begin{itemize}
\item[1)] five different values of $Q(u)$ for $|\eta| > 1$
\item[2)] four different values of $Q(u)$ for $|\eta| = 1$
\item[3)] five different values of $Q(u)$ for $|\eta| < 1$
\item[4)] three different values of $Q(u)$ for $\eta = 0$.
\end{itemize}

%eta = np.linspace(-2,2,400)
%u1 = 1 + eta
%u2 = 1 - eta
%u3 = -1 + eta
%u4 = -1 - eta
%uarr = np.array([u1, u2, u3, u4])
%fig, ax = plt.subplots()
%for i, u in enumerate(uarr):
%    ax.plot(eta, u, label=f"$u_{{{i}}}(\\eta)$")
%ax.legend()

However, the Chern number is zero in the outermost region in each case, since those correspond to a non-zero band gap.
Thus we are really only interested in one, two or at most three values of $Q$.
\\
Generally, the Chern number can be computed using the representation of the Berry curvature as
\begin{align}
\vec{\Omega} &= \frac{\vec{d}}{2d^3},
\end{align}
and integrating over the Brillouin-zone with normal vector $\vec{n} = (0,0,1)^\top$
\begin{align}
Q &= \frac{1}{2\pi}\intx\int_{\tecxt\text{BZ}}^{ }\vec{\Omega}\cdot\mathrm{d}\vec{S} = \frac{1}{2\pi} \intx\int_{\tecxt\text{BZ}}^{ } \frac{d_z}{2d^3}\, \mathrm{d}k_x \tecxt\text{d}k_y = \\
&= \frac{1}{\pi}\intx\int_{0}^{\pi}\intx\int_{0}^{\pi}\frac{u + \eta\cos x + \cos y}{\kla\left( \eta^2\sin^2x + \sin^2y + \klam\left[ u + \eta\cos x + \cos y \right]^2 \right)^{3/2}}\, \mathrm{d}x \mathrm{d}y = \\
&= -\frac{1}{\pi} \partial_u\intx\int_{0}^{\pi}\intx\int_{0}^{\pi} \frac{1}{\sqrt{\eta^2\sin^2x + \sin^2y + \klam\left[ u + \eta\cos x + \cos y \right]^2}} \, \mathrm{d}x \mathrm{d}y = \\
&= -\frac{1}{\pi} \partial_u\intx\int_{0}^{\pi}\intx\int_{0}^{\pi} \frac{1}{\sqrt{\eta^2 + u^2 + 1 + 2u (\eta\cos x + \cos y) + 2\eta\cos x\cos y}} \, \mathrm{d}x \mathrm{d}y = \\
&= ???.
\end{align}
Let $v=\cos x$ and $w=\cos y$.
Then we find
\begin{align}
Q &= -\frac{1}{\pi}\partial_u\intx\int_{-1}^{1}\intx\int_{-1}^{1} \frac{\sqrt{1-v^2}\sqrt{1-w^2}}{\sqrt{\eta^2+u^2+1+2u\eta v + 2uw+2\eta vw}}\, \mathrm{d}v \mathrm{d}w
\end{align}
Obviously this integral is ridiculous, so there must be a more clever way... 
See Gauss-Bonnet theorem for a much deeper theory.
%\begin{align}
%Q &= \frac{1}{2\pi}\intx\int_{-\pi/a}^{\pi/a}\intx\int_{-\pi/a}^{\pi/a} \frac{1}{2d^2}\, \mathrm{d}k_x \mathrm{d}k_y = \\
%&= \frac{1}{\pi}\intx\int_{0}^{\pi}\intx\int_{0}^{\pi} \frac{1}{\eta^2\sin^2x + \sin^2y + \klam\left[ u + \eta\cos x + \cos y \right]^2}\, \mathrm{d}x\, \mathrm{d}y = \\
%&= \frac{1}{\pi}\intx\int_{-1}^{1}\intx\int_{-1}^{1} \frac{\sqrt{1-v^2}\sqrt{1-w^2}}{\eta^2+u^2+1+2u\eta v + 2uw+2\eta vw} \, \mathrm{d}v \mathrm{d}w = \\
%&= \frac{1}{\pi}\intx\int_{-1}^{1}\frac{\sqrt{1-w^2}}{2\eta (u+w)} \intx\int_{-1}^{1} \frac{\sqrt{1-v^2}}{k + v}\, \mathrm{d}v\, \mathrm{d}w = \\
%&= \frac{1}{\pi}\intx\int_{-1}^{1}\frac{\sqrt{1-w^2}}{2\eta (u + w)}\kla\left(  -\frac{{\pi}\sqrt{k^2-1}-{\pi}k}{2}+\arcsin\left(\frac{1}{k}\right)\sqrt{k^2-1}-1 \right) \, \mathrm{d}w
%\end{align}
%Here $k(w) = \frac{\eta^2+u^2+1 + 2uw}{2\eta (u+w)}$.

%def integrand(x, eta=1.0, u=1.0):
%    k = (eta**2 + u**2 + 1 + 2*u*x) / (2*eta * (u + x))
%    fac = -(np.pi*np.sqrt(k**2-1)-np.pi*k)/2 + np.arcsin(1/k)*np.sqrt(k**2-1)-1
%    return np.sqrt(1 - x**2) * fac / (2*eta * (u + x))


%%%https://phyx.readthedocs.io/en/latest/TI/Lecture%20notes/3.html

%@njit
%def chern(u, eta, N=100):
%    kx, ky = np.linspace(-np.pi, np.pi, N+1), np.linspace(-np.pi, np.pi, N+1)
%    res = 0
%    for i in range(N):
%        for j in range(N):
%            dvec = np.array([eta * np.sin(kx[i]), np.sin(ky[j]), u + np.cos(ky[j]) + eta * np.cos(kx[i])])
%            res += dvec[2] / np.linalg.norm(dvec, ord=2)**3
%    return res / (4*np.pi) * (kx[1] - kx[0])**2

%The Chern number may also be obtained by counting how often the image of the Brillouin zone ``wraps around the origin''.
%In the lecture we already discussed the case of $\eta = 1$, which gives $Q(u) = \pm 1$ for $0<u<2$ and $-2<u<0$ respectively.
%But I fail to understand how the ``torus''-like structure came about, which apparently is this BZ-image.

%Figure \ref{figure:phi} shows how the angle $\phi(k)$ behaves for the same parameters. 
%\begin{figure}[ht!]
%\centering
%\includegraphics[width=\linewidth]{01_WindingNumber_phi_v2.png}
%\caption{Function $\phi(k)$ for $k\in[-\pi,\pi)$. }
%\label{figure:phi}
%\end{figure}

%\begin{thebibliography}{9}
%\bibitem{example} description, \url{https://www.exmapleurl}, day.\,month~2021.
%\end{thebibliography}

\end{document}